\chapterAndLabel{Path Shape Checker}{path-shape-checker}

The Path Shape Checker tracks which paths have a name and which have a
parent.  It applies to
\sunjavadoc{java.base/java/nio/file/Path.html}{Path} and
\sunjavadoc{java.base/java/io/File.html}{File} objects.  For example, if
variable \<p> has type \<@HasParent Path>, then \<p.getParent()> returns a
non-\<null> value.

Section~\ref{path-shape-nullness} describes how path shape annotations
enable the Nullness Checker (\chapterpageref{nullness-checker}) to treat
calls to methods such as
\sunjavadoc{java.base/java/nio/file/Path.html\#getParent()}{Path.getParent}
more precisely by refining their result to \<@NonNull> in some cases.
For example, the Nullness Checker accepts this code:

\begin{Verbatim}
  Path parent = Path.of("myfile.txt").toAbsolutePath().getParent();
  parent.toString();  // OK:  parent is non-null
\end{Verbatim}


\sectionAndLabel{Invoking the Path Shape Checker}{path-shape-invoking}

You will not typically run the Path Shape Checker.  The Nullness Checker
runs it automatically.

You can unsoundly suppress warnings related to path shapes with
\<@SuppressWarnings("pathshape")>, or everywhere by using command-line option
\<-AsuppressWarnings=pathshape>; see \chapterpageref{suppressing-warnings}.

The command-line argument \<-AnoPathShape> makes the Nullness Checker
not run the Path Shape Checker.  Then, the Nullness Checker makes no
refinement based on path shapes.


\sectionAndLabel{Path shape annotations}{path-shape-annotations}

These qualifiers are part of the Path Shape type system:

\begin{description}

\item[\refqualclass{checker/nullness/qual}{HasName}]
  indicates that the path has at least one name element and is not the
  empty path.  For a \<Path>, \<getFileName()> returns a non-\<null>,
  non-empty path.  For a \<File>, \<getName()> returns a non-empty string.
  Roots (such as \</> and, on Windows, \<C:\textbackslash{}>, \<C:>, and
  \<\textbackslash{}\textbackslash{}server\textbackslash{}share\textbackslash{}>)
  and the empty path \<""> are not \<@HasName>.

\item[\refqualclass{checker/nullness/qual}{HasParent}]
  indicates that the path has a parent:  \<Path.getParent()>,
  \<File.getParent()>, and \<File.getParentFile()> return a non-\<null>
  value.  A path with a parent has either 2 or more names (such as
  \<a/b>), or a root and a name (such as \</a>).

\item[\refqualclass{checker/nullness/qual}{UnknownPathShape}]
  indicates that the path might be a root or the empty path.
  It is the default type qualifier.

\item[\refqualclass{checker/nullness/qual}{PathShapeBottom}]
  is used internally by the type system but should never be written by a
  programmer.  Its only value is \<null>.

\item[\refqualclass{checker/nullness/qual}{PolyPathShape}]
  is a polymorphic qualifier (Section~\ref{method-qualifier-polymorphism}).

\end{description}

The subtyping hierarchy is a chain:  \<@PathShapeBottom> is a subtype of
\<@HasParent>, which is a subtype of \<@HasName>, which is a subtype of
\<@UnknownPathShape>.

A path's shape depends only on the characters of the path, not on the file
system.  Therefore, once a path is known to have a name or a parent, that
fact remains true even if the file system changes.


\sectionAndLabel{How the Nullness Checker uses path shapes}{path-shape-nullness}

The Nullness Checker refines the result of these methods to \<@NonNull>:

\begin{itemize}
\item
  \<Path.getParent()>, \<File.getParent()>, and \<File.getParentFile()>,
  when the receiver is \<@HasParent>.
\item
  \<Path.getFileName()>, when the receiver is \<@HasName>.
\end{itemize}


\sectionAndLabel{Inferring path shapes}{path-shape-inference}

You rarely need to write path shape annotations, because the Path Shape
Checker infers them.  It uses these rules:

\begin{itemize}
\item
  \<p.toAbsolutePath()> and \<f.getAbsoluteFile()> are \<@HasParent> if
  the receiver is \<@HasName>.  For example, if \<p> is \<a>, then
  \<p.toAbsolutePath()> might be \</home/user/a>.  A path that might be
  empty, such as an unannotated \<Path>, does not have this property, because
  the absolute path of the empty path is the current directory, which might
  be \</>.
\item
  \<p.getFileName()>, \<p.getName(i)>, and \<p.subpath(i, j)> are
  \<@HasName> if the receiver is \<@HasName>.
\item
  \<a.resolve(b)> is \<@HasParent> if \<b> is \<@HasParent> or if both
  \<a> and \<b> are \<@HasName>.  Otherwise, it is \<@HasName> if \<b> is
  \<@HasName>.  \<a.resolveSibling(b)> is \<@HasParent> if \<b> is
  \<@HasParent> or if \<a> is \<@HasParent> and \<b> is \<@HasName>.
  Otherwise, it is \<@HasName> if \<b> is \<@HasName>.
\item
  \<p.toFile()> and \<f.toPath()> have the same path shape as the receiver.
  (On Windows, \<f.toPath()> does not always preserve the path shape; see
  Section~\ref{path-shape-limitations}.)
\item
  \<Path.of>, \<Paths.get>, and \<new File(String)>, and the \<String>
  overloads of \<resolve> and \<resolveSibling>, have a path shape that is
  determined from the string, if every string argument is a compile-time
  constant.  The Path Shape Checker determines the path shape under both
  Unix and Windows rules and uses the weaker result.  For example, \<"C:">
  is a name on Unix but a root on Windows, so \<Path.of("C:")> is
  \<@UnknownPathShape>.  The Path Shape Checker does not examine strings that
  are not compile-time constants.
\item
  \<new File(parent, child)> is \<@HasParent> if \<child> is a compile-time
  constant with a parent, such as \<"a/b">, and \<@HasName> if \<child> is a
  compile-time constant with a name, such as \<"a">.
\item
  If \<Files.isRegularFile(p)> or \<f.isFile()> returns true, then \<p> or
  \<f> is \<@HasName>, because a root and the empty path always denote a
  directory.
\item
  Elements of \<Files.list(dir)>, \<Files.newDirectoryStream(dir)>, and
  \<dir.listFiles()>, and the results of \<Files.createTempFile(dir, ...)>
  and \<Files.createTempDirectory(dir, ...)>, are \<@HasParent> if \<dir>
  is \<@HasName>, and \<@HasName> otherwise.  The results of
  \<Files.createTempFile(prefix, suffix)> and
  \<Files.createTempDirectory(prefix)> are \<@HasName>.
\item
  Elements of \<Files.walk(start)> and \<Files.find(start, ...)> have the
  same path shape as \<start>, because the first element is \<start>
  itself.
\item
  \<p.normalize()>, \<p.toRealPath()>, and \<f.getCanonicalFile()> are
  \<@UnknownPathShape>.  For example, \<Path.of("a/..").normalize()> is the
  empty path.
\end{itemize}

Some of these rules qualify type arguments, as in \<Stream<@HasName Path>>.
Ordinarily, type arguments are invariant
(Section~\ref{covariant-type-parameters}).  The Path Shape Checker treats the
type argument of \<Iterable>, \<Iterator>, \<Spliterator>, \<Optional>,
\<Stream>, and \<DirectoryStream> as covariant, so it accepts code such
as \<Stream<Path> s = Files.list(dir);>.  This is sound, because these
interfaces have no methods that insert elements.
Likewise, \<Stream.toList()> returns an unmodifiable list, so the Path Shape
Checker treats the type argument of its result as covariant.  It accepts
code such as \<return Files.list(dir).toList();> in a method whose return
type is \<List<Path>>.

An effectively final local variable with an initializer takes the covariant
type arguments of its initializer, as if the variable had been declared with
\<var>.  This changes only a type argument whose qualifier is
\<@UnknownPathShape>, whether that qualifier is written or defaulted.  For
example, the Path Shape Checker accepts this code:

\begin{Verbatim}
  try (Stream<Path> s = Files.list(Path.of("mydir"))) {
    s.forEach(p -> p.getParent().toString());   // OK:  p is @HasParent
  }
\end{Verbatim}

\noindent
The variable does not take other type arguments, such as that of \<List>,
nor the component type of an array.  Doing so would make the Path Shape
Checker reject correct code that uses the variable where its declared type
is expected, such as passing a \<List<Path>> variable to a method whose
formal parameter type is \<List<Path>>, or storing a \<File> whose path
shape is unknown into a \<File[]> variable.  To make a variable take all
the type arguments of its initializer, declare it with \<var>.


\sectionAndLabel{Limitations}{path-shape-limitations}

The Path Shape Checker assumes that implementations of \<Path> follow the
specification of \<Path>.  \<Path> is an interface, so a third-party
\sunjavadoc{java.base/java/nio/file/spi/FileSystemProvider.html}{FileSystemProvider}
could violate that assumption.

The Path Shape Checker assumes that \<File.toPath()> preserves the path
shape.  That is not true on Windows:  \<new File("//server/share")> has a
parent, \<\textbackslash{}\textbackslash{}server>, but its \<toPath()> is a
root.  To make the Path Shape Checker sound for such paths, pass the
command-line argument \<-Astubs=sound-file-topath.astub> to the Nullness
Checker.  Then, \<File.toPath()> returns \<@UnknownPathShape>.

A rule that refines a type argument can make the Path Shape Checker reject
some correct code.  For example, \<Files.list(dir).findFirst()> has type
\<Optional<@HasName Path>>, so the Path Shape Checker rejects
\<Files.list(dir).findFirst().orElse(dir)> unless \<dir> is
\<@HasName>.  The same is true if the result of \<findFirst()> is the
initializer of an effectively final local variable of type
\<Optional<Path>>, because the variable takes the covariant type arguments of
its initializer.  You can cast the result of \<findFirst()> to
\<Optional<Path>>, or you can suppress the warning.

%% LocalWords:  pathshape HasName HasParent UnknownPathShape PathShapeBottom
%% LocalWords:  PolyPathShape AnoPathShape getParent getParentFile getName
%% LocalWords:  getFileName toAbsolutePath getAbsoluteFile resolveSibling
%% LocalWords:  toFile toPath isRegularFile isFile listFiles subpath dir
%% LocalWords:  newDirectoryStream createTempFile createTempDirectory toList
%% LocalWords:  toRealPath getCanonicalFile Spliterator DirectoryStream
%% LocalWords:  FileSystemProvider AsuppressWarnings myfile
